\documentclass{beamer}

\usetheme{kth}
%\usetheme[title style=centered]{kth}
%\usefonttheme{kth}

\title[Short Paper Title]
{Title As It Is In the Proceedings}

\subtitle
{Include Only If Paper Has a Subtitle}

\author[Author, Another] % (optional, use only with lots of authors)
{F.~Author\inst{1} \and S.~Another\inst{2}}

\institute[Universities of Somewhere and Elsewhere] % (optional, but mostly needed)
{
  \inst{1}%
  University of Somewhere
  \qquad
  \inst{2}%
  University of Elsewhere}

\date[CFP 2003] % (optional, should be abbreviation of conference name)
{2019-12-02}

\kthfooterstyle{lcr}
%\kthlfoot{Left}   % Default is <DATE>
%\kthcfoot{Center} % Default is <TITLE>
%\kthrfoot{Right}  % Default is <SLIDE NUMBER>/<TOTAL SLIDES>
\kthlfoot{\today}
\kthcfoot{}

\AtBeginSubsection[]
{
  \begin{frame}<beamer>{Outline}
    \tableofcontents[currentsection,currentsubsection]
  \end{frame}
}
\begin{document}


\frame[plain, t]{\titlepage}

\begin{frame}{Outline}
  \tableofcontents
\end{frame}

\section{Motivation}

\subsection{The Basic Problem That We Studied}

\begin{frame}{Make Titles Informative. Use Uppercase.}{Subtitles are optional.}
  \begin{itemize}
  \item
    Use \texttt{itemize} a lot.
  \item
    Use very short sentences or short phrases.
  \end{itemize}
\end{frame}

\begin{frame}{Make Titles Informative.}
  You can create overlays\dots
  \begin{itemize}
  \item using the \texttt{pause} command:
    \begin{itemize}
    \item
      First item.
      \pause
    \item
      Second item.
    \end{itemize}
  \item
    using overlay specifications:
    \begin{itemize}
    \item<3->
      First item.
    \item<4->
      Second item.
    \end{itemize}
  \item
    using the general \texttt{uncover} command:
    \begin{itemize}
      \uncover<5->{\item
        First item.}
      \uncover<6->{\item
        Second item.}
    \end{itemize}
  \end{itemize}
\end{frame}

\subsection{Previous Work}

\begin{frame}{Make Titles Informative.}
  \begin{columns}
    \column{0.45\textwidth}
      Left column with items
      \begin{itemize}
      \item
        An item to the left
      \item
        Another item to the left
      \item
        ...
      \end{itemize}
    \column{0.45\textwidth}
      \begin{figure}
        \includegraphics[width=0.75\textwidth]{img/logo.pdf}
        \caption{Relative \texttt{\symbol{92}textwidth}}
      \end{figure}
  \end{columns}
\end{frame}

\begin{frame}{Make Titles Informative.}
  \begin{columns}[T]
    \column{0.45\textwidth}
      Left column with items
      \begin{itemize}
      \item
        Use the \texttt{[T]} option to align columns vertically
      \item
        Right column with two images
      \item
        \texttt{\symbol{92}minipage} on the right to split vertically
      \end{itemize}
    \column{0.45\textwidth}
      \begin{minipage}[c][0.35\textheight][c]{0.9\textwidth}
        \begin{figure}
          \includegraphics[width=0.8\textwidth]{img/logo.pdf}
        \end{figure}
      \end{minipage}
      \begin{minipage}[c][0.35\textheight][c]{0.9\textwidth}
        \begin{figure}
          \includegraphics[width=0.4\textwidth]{img/logo.pdf}%
          \includegraphics[width=0.4\textwidth]{img/logo.pdf}
        \end{figure}
      \end{minipage}
  \end{columns}
\end{frame}

\section{Our Results/Contribution}

\subsection{Main Results}

\begin{frame}{Make Titles Informative.}
\end{frame}

\subsection{Basic Ideas for Proofs/Implementation}

\begin{frame}{Make Titles Informative.}
\end{frame}

\section*{Summary}

\begin{frame}{Summary}
  \begin{itemize}
  \item
    The \alert{first main message} of your talk in one or two lines.
  \item
    The \alert{second main message} of your talk in one or two lines.
  \item
    Perhaps a \alert{third message}, but not more than that.
  \end{itemize}

  \vskip0pt plus.5fill
  \begin{itemize}
  \item
    Outlook
    \begin{itemize}
    \item
      Something you haven't solved.
    \item
      Something else you haven't solved.
    \end{itemize}
  \end{itemize}
\end{frame}

\appendix
\section<presentation>*{\appendixname}
\subsection<presentation>*{For Further Reading}

\begin{frame}[allowframebreaks]
  \frametitle<presentation>{For Further Reading}

  \begin{thebibliography}{10}

  \beamertemplatebookbibitems

  \bibitem{Author1990}
    A.~Author.
    \newblock {\em Handbook of Everything}.
    \newblock Some Press, 1990.


  \beamertemplatearticlebibitems

  \bibitem{Someone2000}
    S.~Someone.
    \newblock On this and that.
    \newblock {\em Journal of This and That}, 2(1):50--100,
    2000.
  \end{thebibliography}
\end{frame}

\end{document}


